% Generated draft: requires mathematical and source review.
\documentclass[11pt]{article}
\usepackage[margin=25mm]{geometry}
\usepackage{fontspec}
\setmainfont{Latin Modern Roman}
\usepackage{amsmath,amssymb,mathtools,unicode-math}
\setmathfont{Latin Modern Math}
\usepackage{enumitem,needspace,xurl,hyperref}
\setlength{\parindent}{0pt}
\setlength{\parskip}{5pt}
\setlength{\emergencystretch}{3em}
\newcommand{\N}{\mathbb{N}}
\newcommand{\Z}{\mathbb{Z}}
\newcommand{\Q}{\mathbb{Q}}
\newcommand{\R}{\mathbb{R}}
\newcommand{\C}{\mathbb{C}}
\newcommand{\F}{\mathbb{F}}
\newcommand{\A}{\mathbb{A}}
\newcommand{\PP}{\mathbb{P}}
\newcommand{\E}{\mathbb{E}}
\newcommand{\eps}{\varepsilon}
\newcommand{\dd}{\,\mathrm{d}}
\begin{document}
\section*{UnsolvedMath ID: 30006488}
\textbf{Source problem code:} catalogue entry
\subsection{Definitions and mathematical statement}
\emph{Coefficient qualification.} The supplied record does not specify the coefficient domain or characteristic. The following is a field-based formulation for a fixed field \(k\); whether the intended questions concern every field or a restricted class of fields remains unresolved.

A ring is left Noetherian if every ascending chain of left ideals stabilizes, and right Noetherian is defined analogously. For the enveloping algebras and group algebras below, the standard maps induced by \(x\mapsto -x\) and \(g\mapsto g^{-1}\), respectively, identify each algebra with its opposite algebra, so left and right Noetherianity are equivalent. Here \emph{Noetherian} means this equivalent condition.

For a Lie algebra \(\mathfrak g\) over \(k\), its universal enveloping algebra is
\[
U_k(\mathfrak g)=T_k(\mathfrak g)\big/\left\langle x\otimes y-y\otimes x-[x,y]:x,y\in\mathfrak g\right\rangle,
\]
where \(T_k(\mathfrak g)\) is the tensor algebra and the brackets generate a two-sided ideal. The first question is whether
\[
\text{for every Lie algebra }\mathfrak g\text{ over }k,\qquad
U_k(\mathfrak g)\text{ Noetherian}\ \Longrightarrow\ \dim_k\mathfrak g<\infty.
\]

For a group \(G\), the group algebra \(kG\) consists of finite formal sums \(\sum_{g\in G}a_g g\), with coefficients in \(k\), and multiplication extends the group operation bilinearly. A group \(P\) is polycyclic if it has a finite subnormal series
\[
\{1\}=P_0\triangleleft P_1\triangleleft\cdots\triangleleft P_r=P
\]
whose factors \(P_i/P_{i-1}\) are cyclic, either finite or infinite. A group \(G\) is polycyclic-by-finite if it has a polycyclic normal subgroup of finite index. The second question is whether
\[
\text{for every group }G,\qquad
kG\text{ Noetherian}\ \Longrightarrow\ G\text{ polycyclic-by-finite}.
\]
No finite generation assumption on \(\mathfrak g\) or \(G\) is supplied. These are two separate questions in the extracted statement. The retained original passage supplies background about classification of Hopf algebras of finite Gelfand--Kirillov dimension and difficulties for Noetherian Hopf algebras, but stops at ``Conjecture 1'' without reproducing its mathematical content. That background does not explicitly impose a finite Gelfand--Kirillov dimension hypothesis on either question.
\subsection{Short English statement}
For a fixed coefficient field, ask whether a Noetherian universal enveloping algebra forces its Lie algebra to be finite-dimensional, and whether a Noetherian group algebra forces its group to have a polycyclic normal subgroup of finite index. The intended coefficient fields and characteristic restrictions require confirmation. The supplied record labels the problem open; this draft does not independently verify that status.
\subsection{Sources}
\begin{itemize}
\item Supplied local UnsolvedMath registry record, permanent ID \(30006488\), source problem code \texttt{catalogue entry}, in the Oberwolfach Reports Open Problems collection. No individual proposer is identified; the supplied proposed year is 2025.
\item Supplied report citation: \emph{Mini-Workshop: The Yang--Baxter Equation and Representations of Braid Groups (2026)}, Oberwolfach Reports, pp.~2633--2672, DOI \texttt{10.4171/owr/2025/49}. \url{https://doi.org/10.4171/owr/2025/49}. The literature-source label is ``Oberwolfach Report 49/2025.'' The citation's 2026 designation and its 2025 report designation are retained without reconciliation.
\item Supplied literature source: \emph{Universal enveloping algebra of the Witt algebra}. \url{https://doi.org/10.1016/j.aim.2014.05.007}. No authors or further bibliographic details are supplied.
\item Literature assessment dated exactly \texttt{2026-08-22}: the supplied curation labels the problem open and describes the Noetherian enveloping-algebra conjecture as a difficult classical open problem, stating that known non-Noetherian examples and graded results do not settle it. This is supplied dated curation derived from the OWR record, not an independently verified literature assessment. Time-sensitive status and cited claims require independent verification; the assessment does not separately assess the group-algebra question.
\item The registry labels the statement \texttt{corrected\_verified} and reports that the official Oberwolfach report was checked, the question isolated from surrounding exposition, and notation repaired. These are supplied source-verification claims, not verification performed for this draft. The registry's open label likewise remains a supplied classification.
\end{itemize}
\end{document}
